\documentclass[11pt]{article}
\usepackage[margin=1in]{geometry}
\usepackage[T1]{fontenc}
\usepackage{lmodern,microtype,amsmath,amssymb,amsthm,mathtools}
\usepackage{booktabs,array,longtable,enumitem}
\usepackage{xcolor}
\usepackage[colorlinks=true,linkcolor=blue!50!black,citecolor=blue!50!black,urlcolor=blue!50!black]{hyperref}
\hypersetup{pdftitle={Conservation laws and exterior flux for heat-zero hierarchies},pdfauthor={Edward Baker}}
\newtheorem{theorem}{Theorem}[section]
\newtheorem{proposition}[theorem]{Proposition}
\newtheorem{lemma}[theorem]{Lemma}
\newtheorem{corollary}[theorem]{Corollary}
\theoremstyle{remark}
\newtheorem{remark}[theorem]{Remark}
\newcommand{\dd}{\,\mathrm d}
\newcommand{\ii}{\mathrm i}
\newcommand{\C}{\mathbb C}
\newcommand{\R}{\mathbb R}
\newcommand{\av}[1]{\frac{1}{2\pi\ii}\oint_\Gamma #1\dd z}
\newcommand{\F}{\mathcal F}
\newcommand{\E}{\mathcal E}
\DeclareMathOperator{\tr}{tr}
\DeclareMathOperator{\spec}{spec}
\numberwithin{equation}{section}
\setcounter{tocdepth}{1}
\setlist{itemsep=3pt,topsep=5pt}
\title{Conservation laws and exterior flux\protect\\for heat-zero hierarchies}
\author{Edward Baker\thanks{Draft prepared for Edward Baker with substantial OpenAI Codex LLM assistance. Actual model and reasoning effort were unavailable in the guarded conversation metadata; none are inferred.}}
\date{18 November 2026\\\small Working draft; same-assistant checks, no independent review}
\begin{document}
\maketitle
\begin{abstract}
We investigate conserved quantities in the exterior-moment hierarchy of
the physical Newman heat equation $H_t=-H_{zz}$. Finite polynomial heat
flows admit explicit coefficient integrals and autonomous symmetric
integrals, including a centered fourth-order statistic. We derive these
integrals, a restricted obstruction to universal finite-moment integrals,
and the distinction between full-field heat pullback and a useful local
estimate. For a counted block in a larger divisor, every homogeneous,
translation-invariant polynomial heat integral has an exact exterior-flux
law. Removing its affine exterior response yields one Cauchy kernel whose
first two Laurent terms cancel. On an enclosing circle this gives a
quadratic radius-factor reduction of a specified separate source-error
allowance; a centered symmetric quartet has a cubic refinement. We also
pay uncertain block coordinates explicitly. Exact finite examples retain
both nonzero exterior flux and a collision with positive fourth-order
invariant. The results supply identities and intermediate error estimates,
not a canonical gap improvement or a collision-exclusion theorem.
\end{abstract}
\tableofcontents
\clearpage
\section{Physical equation, provenance, and the pair balance}
\label{sec:context}

The parent manuscript \cite{BakerExterior2026} develops a complete
exterior-moment hierarchy and counted-block transport for the Newman
field. The base manuscript \cite{BakerVolterra2026} supplies its canonical
source and Jacobi--Volterra realization. We use their normalization
throughout:
\begin{align}
 r_k(u)&=\pi k^2e^{4u},&
 \Phi(u)&=\sum_{k\ge1} e^u\bigl(2r_k(u)^2-3r_k(u)\bigr)e^{-r_k(u)},
 \label{eq:source}\\
 H(t,z)&=\int_0^\infty e^{tu^2}\Phi(u)\cos(zu)\dd u,&
 H_t&=-H_{zz},\qquad H(0,z)=\tfrac18\xi((1+\ii z)/2).
 \label{eq:heat}
\end{align}
For $u\ge0$, the source is positive and decreases faster than
$e^{-Cu^2-Cu}$ for every fixed $C>0$. Compact differentiation of
\eqref{eq:heat} is therefore legitimate to every fixed order. For global
divisor statements we use the parent's established even order-one
canonical product, including its count, cutoff and complete
regularization hypotheses. The local block results below need only a
holomorphic heat solution and a fixed zero-free enclosing boundary.
They do not require an all-real exterior.

The founding exploration \cite{BakerConservationBrainstorm2026} contains
the finite identities and fourth-order block balance developed here.
The bounded extension in Sections~\ref{sec:block}--\ref{sec:errors}
selects a general homogeneous class and compares explicit error
allowances. The result is an identity and intermediate-estimate
investigation in the terminology of the project's result and avoidance
indices \cite{ProjectResults2026,ProjectAvoid2026}. A representation of
initial data does not by itself supply a bound on that data. A negative
result is always restricted to its displayed class. We make no priority
claim for finite heat-pole integrability: the abstract of
\cite{DeconinckKimuraSegur2007} reports conserved quantities for rational
Burgers poles, and the checked passages of \cite{Sokal2023HeatZeros}
discuss finite heat-zero dynamics and the specialized matrix in
Section~\ref{sec:calogero}. All formulas used here are derived below;
no uninspected external theorem or conjecture is imported.

\subsection{Count and pair data}
Put $B=H_z/H$. Off the divisor,
\begin{equation}
 B_t=-\partial_z(B_z+B^2).
 \label{eq:burgers}
\end{equation}
If one stationary positively oriented contour $\Gamma$ remains zero-free,
its integer count
\begin{equation}
 N=\av{B(t,z)}
 \label{eq:count}
\end{equation}
is constant. This follows either from continuity and the argument
principle, or by integrating the total derivative in
\eqref{eq:burgers}. Two simple zeros becoming one double zero do not
change the multiplicity count.

For a counted pair write its roots as $c\pm s$, set $D=s^2$, and let
$w_j=\rho_j-c$ for each exterior occurrence. The parent's complete
collar assumptions make the following sums absolutely convergent:
\begin{equation}
 S_n=\sum_j(w_j^2-D)^{-n}\quad(n\ge1),\qquad
 T_n=\sum_jw_j(w_j^2-D)^{-n}\quad(n\ge2).
 \label{eq:moments}
\end{equation}
In particular, $T_1$ is not part of this convergent coordinate family.
Writing $M=S_1$ and $\F=12S_2+16DS_3$, the regular pair equations are
\begin{equation}
 D'=2-4DM,\qquad M'=2M^2-\F.
 \label{eq:pair}
\end{equation}
The first law follows by subtracting the two simple-root velocities;
the mutual force contributes $2$ to $D'$, and each exterior root
contributes $-4D/(w_j^2-D)$. The second is the $n=1$ row of the
integer hierarchy displayed in Section~\ref{sec:obstruction}. The
parent's counted-contour passage extends both laws through mergers.

\begin{proposition}[An energy balance, not a general first integral]
Define
\begin{equation}
 E=2M-2DM^2+D\F.
 \label{eq:pair-energy}
\end{equation}
Then $E'=D\F'$. Also,
\begin{equation}
 (D')^2-2DD''+8\F D^2=4.
 \label{eq:pair-identity}
\end{equation}
For a prescribed forcing $\F(t)$, \eqref{eq:pair} is the
time-dependent Hamiltonian system $D'=E_M$, $M'=-E_D$, where partial
derivatives hold $\F(t)$ fixed.
\end{proposition}
\begin{proof}
Substitute \eqref{eq:pair} into
\[
 E'=2M'-2D'M^2-4DMM'+D'\F+D\F'.
\]
All terms except $D\F'$ cancel. Differentiating the first equation of
\eqref{eq:pair} gives $D''=-8M+8DM^2+4D\F$, whose substitution
proves \eqref{eq:pair-identity}. The two partial derivatives of
\eqref{eq:pair-energy} give the asserted Hamilton equations.
\end{proof}

The physical hierarchy determines $\F'$ through higher moments, so
constant forcing cannot be assumed. Subtracting $\int D\F'\dd t$
formally produces a constant but still requires that entire flux.
At a pair merger $D=0$, \eqref{eq:pair} gives $D'=2$, consistently
with \eqref{eq:pair-identity}. On a chart where $D>0$, putting
$r=\sqrt D$ gives
\begin{equation}
 r''=2\F r-r^{-3}.
\end{equation}
The inverse-cubic term is attractive. These formulas do not produce
a positive barrier against a merger; the $r$ chart itself ends there.

\section{Complete finite heat divisors}
\label{sec:finite}

Let $P(t,z)=\sum_{j=0}^N a_j(t)z^{N-j}$ be monic, with $a_0=1$,
and suppose $P_t=-P_{zz}$. This section concerns the complete finite
polynomial, with no omitted exterior. We allow complex coefficients
and roots unless a reality condition is explicitly stated.

\begin{theorem}[Coefficient integrals and a polynomial clock]
\label{thm:finite}
The coefficients satisfy
\begin{equation}
 a_j'=-(N-j+2)(N-j+1)a_{j-2},\qquad a_{-1}=a_{-2}=0.
 \label{eq:coeff-law}
\end{equation}
Consequently
\begin{equation}
 I_j(t)=\sum_{r=0}^{\lfloor j/2\rfloor}
 \frac{t^r}{r!}\frac{(N-j+2r)!}{(N-j)!}a_{j-2r}(t)=a_j(0).
 \label{eq:explicit-integrals}
\end{equation}
For $N\ge2$, set $\theta=-a_2/[N(N-1)]$. Then $\theta'=1$ and
every coefficient of
\begin{equation}
 Z(z)=e^{\theta\partial_z^2}P(t,z)
 \label{eq:autonomous}
\end{equation}
is constant. Its coefficient of $z^{N-2}$ is zero. Its other $N-1$
nonleading coefficients are independent autonomous integrals of the
$N$-dimensional monic coefficient flow.
\end{theorem}
\begin{proof}
Comparing powers in $P_t=-P_{zz}$ gives \eqref{eq:coeff-law}.
The exponential of $\partial_z^2$ terminates on a polynomial and
commutes with differentiation in time. Thus
$\partial_t(e^{t\partial_z^2}P)=0$. Its $j$th coefficient is
\eqref{eq:explicit-integrals}. Similarly,
\[
 \partial_t(e^{\theta\partial_z^2}P)
 =(\theta'-1)\partial_z^2e^{\theta\partial_z^2}P=0.
\]
The $z^{N-2}$ coefficient of $Z$ is $a_2+N(N-1)\theta=0$.
Conversely, an arbitrary monic $Z$ with that coefficient zero,
together with an arbitrary $\theta$, recovers
$P=e^{-\theta\partial_z^2}Z$. Its second coefficient is
$-N(N-1)\theta$, so this is an inverse polynomial coordinate change.
Independence and completeness of the $N-1$ constants follow.
\end{proof}

For simple roots, evaluation of the heat equation at $P(t,x_i(t))=0$
gives
\begin{equation}
 x_i'=\frac{P_{zz}(t,x_i)}{P_z(t,x_i)}
       =2\sum_{j\ne i}\frac1{x_i-x_j}.
 \label{eq:finite-roots}
\end{equation}
Set
\begin{equation}
 c=\frac1N\sum_i x_i,\qquad y_i=x_i-c,\qquad p_k=\sum_i y_i^k.
 \label{eq:center}
\end{equation}
Then $p_0=N$, $p_1=0$, and pairwise cancellation gives $c'=0$.
For any integer $k\ge2$,
\begin{equation}
 p_k'=2k\sum_{i<j}\frac{y_i^{k-1}-y_j^{k-1}}{y_i-y_j}.
 \label{eq:powers}
\end{equation}
In particular,
\begin{equation}
 p_2'=2N(N-1),\qquad p_3'=0,\qquad p_4'=4(2N-3)p_2.
 \label{eq:low-powers}
\end{equation}
For the last identity, the unordered sum of $y_i^2+y_iy_j+y_j^2$
is $(N-1)p_2-p_2/2=(N-3/2)p_2$. This proves every coefficient
in \eqref{eq:low-powers}, including its sign.

\begin{corollary}[Centered homogeneous integrals]
\label{cor:homogeneous}
Let $Q(t,y)=\prod_i(y-y_i)=\sum_{j=0}^N b_jy^{N-j}$, so $b_1=0$
and $b_2=-p_2/2$. Define
\begin{equation}
 \vartheta=\frac{p_2}{2N(N-1)},\qquad
 J_j=[y^{N-j}]\,e^{\vartheta\partial_y^2}Q,
 \quad 3\le j\le N.
 \label{eq:family}
\end{equation}
Each $J_j$ is a translation-invariant symmetric polynomial in the
$x_i$, homogeneous of degree $j$, and a first integral of
\eqref{eq:finite-roots}. The $N-2$ polynomials $J_3,\ldots,J_N$
are independent in the centered coefficient variables. Along with
the centroid they give $N-1$ independent autonomous integrals.
\end{corollary}
\begin{proof}
The center is constant, so $Q_t=-Q_{yy}$ and $\vartheta'=1$.
Theorem~\ref{thm:finite} applies to $Q$, and its inverse coordinate
change now uses $\vartheta,J_3,\ldots,J_N$. Translation leaves every
$y_i$ unchanged. Under $x\mapsto\lambda x$, the centered coefficient
$b_{j-2r}$ has degree $j-2r$ and $\vartheta^r$ has degree $2r$.
Every term in $J_j$ therefore has degree $j$. The coefficient formula
also proves polynomiality and symmetry without selecting root labels.
\end{proof}

Newton identities give $J_3=b_3=-p_3/3$. For $N\ge4$ they give
the useful fourth-order integral
\begin{align}
 C_4:=J_4
 &=\frac{2N-3}{4N(N-1)}p_2^2-\frac14p_4 \label{eq:C4}\\
 &=b_4-\frac{(N-2)(N-3)}{2N(N-1)}b_2^2.
 \notag
\end{align}
Indeed $b_4=p_2^2/8-p_4/4$ and $b_2=-p_2/2$; substitution into
\eqref{eq:family} proves \eqref{eq:C4}. It can also be verified
directly from \eqref{eq:low-powers}. All these statements extend
through collisions because they are polynomial coefficient identities.

\subsection{A collision with positive conserved value}
For $b>0$ and arbitrary real $T$, consider the exact heat polynomial
\begin{equation}
 P(t,z)=(z^2-b^2)^2-(t-T)(12z^2-4b^2)+12(t-T)^2.
 \label{eq:quartic-control}
\end{equation}
Writing $\tau=t-T$, its coefficients are
$a=-2b^2-12\tau$ and $d=b^4+4b^2\tau+12\tau^2$ in
$P=z^4+az^2+d$. Direct differentiation verifies $P_t=-P_{zz}$,
and \eqref{eq:C4} becomes
\begin{equation}
 C_4=d-a^2/12=2b^4/3>0.
 \label{eq:positive-collision}
\end{equation}
Nevertheless $\pm b$ are double zeros at $t=T$.

More generally, for a real centered even quartic $z^4+az^2+d$,
$p_2=-2a$ and $d=C_4+p_2^2/48$. The discriminant in $z^2$ is
\begin{equation}
 \Delta_{z^2}=a^2-4d=p_2^2/6-4C_4.
 \label{eq:quartic-discriminant}
\end{equation}
If $p_2>0$ and $C_4>0$, both $z^2$ roots have positive sum and
positive product. Thus the four $z$ roots are real and simple exactly
when $p_2^2>24C_4$. Equality gives the two nonzero double roots.
This finite threshold uses the variance clock and the invariant jointly.
It contains exactly the information of the direct quartic discriminant;
no same-information advantage is asserted. Positivity of $C_4$ alone
does not exclude a collision.

\section{The specialized finite Calogero matrix}
\label{sec:calogero}

The first-order law \eqref{eq:finite-roots} embeds into a second-order
system, but its usual spectral quantities have a special value on this
locus. This is a finite-dimensional assertion for distinct roots.

\begin{proposition}[Acceleration and zero energy]
If $s_i=\sum_{j\ne i}(x_i-x_j)^{-1}$ and $x_i'=2s_i$, then
\begin{equation}
 x_i''=-8\sum_{j\ne i}(x_i-x_j)^{-3},\qquad
 \mathcal H=\frac12\sum_i(x_i')^2
             -4\sum_{i<j}(x_i-x_j)^{-2}=0.
 \label{eq:calogero-energy}
\end{equation}
\end{proposition}
\begin{proof}
Put $d_{ij}=x_i-x_j$. Then
\[
 s_i-s_j=\frac2{d_{ij}}
          -d_{ij}\sum_{k\ne i,j}\frac1{d_{ik}d_{jk}}.
\]
Differentiating the velocity gives
\[
 x_i''=-4\sum_{j\ne i}\frac{s_i-s_j}{d_{ij}^2}
       =-8\sum_{j\ne i}d_{ij}^{-3}
         +4\sum_{\substack{j,k\ne i\\j\ne k}}
                      \frac1{d_{ij}d_{ik}d_{jk}}.
\]
The last double sum cancels under interchange of $j$ and $k$.
In the expansion of $\sum_i s_i^2$, the squared terms contribute
$2\sum_{i<j}d_{ij}^{-2}$. For each unordered triple of distinct
nodes, the remaining contribution is twice
\[
 \frac1{(x_i-x_j)(x_i-x_k)}
 +\frac1{(x_j-x_i)(x_j-x_k)}
 +\frac1{(x_k-x_i)(x_k-x_j)}=0.
\]
Thus $\sum_i s_i^2=2\sum_{i<j}d_{ij}^{-2}$, proving the energy
identity. For complex roots these are algebraic squares, not absolute
squares; the identity does not assert positive kinetic energy.
\end{proof}

\begin{theorem}[Nilpotence on the heat-flow locus]
\label{thm:nilpotence}
For any $N$ distinct complex nodes, define
\begin{equation}
 L_{ii}=2s_i,\qquad L_{ij}=\frac2{x_i-x_j}\quad(i\ne j).
 \label{eq:L}
\end{equation}
Then $L^N=0$ and $\tr L^k=0$ for every $k\ge1$.
\end{theorem}
\begin{proof}
Let $Q(z)=\prod_i(z-x_i)$ and identify polynomials of degree less
than $N$ with their values at the nodes. The cardinal polynomial at
$x_j$ is $\ell_j(z)=Q(z)/[(z-x_j)Q'(x_j)]$. The matrix $A$ of
differentiation in these value coordinates is therefore
\[
 A_{ii}=\ell_i'(x_i)=\frac{Q''(x_i)}{2Q'(x_i)}=s_i,
 \qquad
 A_{ij}=\ell_j'(x_i)=\frac{Q'(x_i)}{Q'(x_j)(x_i-x_j)}.
\]
Differentiating $N$ times annihilates every such polynomial, so
$A^N=0$. With $W=\operatorname{diag}(Q'(x_i))$, which is invertible,
$L=2W^{-1}AW$. Nilpotence and the vanishing traces follow.
\end{proof}

This proof gives the specialized matrix statement directly; related
finite matrices appear on the checked slides 27--28 of
\cite{Sokal2023HeatZeros}. No conjectural assertion from those slides is
needed. Nilpotence does not eliminate the nonzero coefficient integrals
of Theorem~\ref{thm:finite}; it says that these particular spectral
powers are all zero. It also says nothing about a regularized trace of
an infinite operator. Such a trace would need its own domain,
convergence, differentiation and boundary-flux construction.

\section{A restricted obstruction for finite moment heads}
\label{sec:obstruction}

The integer hierarchy of \cite[Section 2, Integer moment evolution]{BakerExterior2026}
is
\begin{align}
 S_n'={}&2n\sum_{k=1}^n S_kS_{n+1-k}
    +2n\sum_{k=2}^n\bigl(T_kT_{n+2-k}+DS_kS_{n+2-k}\bigr)
    \notag\\
    &-4n(n+2)S_{n+1}-4n(n+3)DS_{n+2},\qquad n\ge1,
    \label{eq:Shierarchy}\\
 T_n'={}&2(2n-1)\sum_{k=2}^nT_kS_{n+1-k}
    +4nD\sum_{k=2}^nT_kS_{n+2-k}
    \notag\\
    &-4(n^2+n-1)T_{n+1}-4n(n+3)DT_{n+2},\qquad n\ge2.
    \label{eq:Thierarchy}
\end{align}
Empty sums vanish. The nonlinear terms in the $n$th row use indices
at most $n$; the two higher indices occur linearly. This feature, with
the pair law, is the only input needed for the obstruction below.

For clarity about the inherited physical realization, at a finite
simple polynomial slice one may temporarily use
$P_1=\sum_j w_j/(w_j^2-D)$. The exterior velocities relative to the
pair mean are
\begin{equation}
 w_j'=4w_jm_j+2P_1+2\sum_{k\ne j}(w_j-w_k)^{-1},
 \qquad m_j=(w_j^2-D)^{-1}.
 \label{eq:relative}
\end{equation}
Differentiate the weights in \eqref{eq:moments}, pair the last sums
using divided differences, and subtract their diagonals. If
$A=(u^2-D)^{-1}$, $C=(v^2-D)^{-1}$ and
$L_h=\sum_{r=0}^{h-1}A^{h-1-r}C^r$ with $L_0=0$, the identities are
\begin{equation}
 \frac{A^h-C^h}{u-v}=-(u+v)AC L_h,\qquad
 \frac{uA^h-vC^h}{u-v}=-AC L_{h-1}-(D+uv)AC L_h.
 \label{eq:divdiff}
\end{equation}
The endpoints of the convolutions cancel every occurrence of $P_1$,
leaving \eqref{eq:Shierarchy}--\eqref{eq:Thierarchy}. These rational
divided differences also remain regular when two exterior nodes
coincide away from the target collar. The parent's proof supplies the
complete canonical passage using cleared-contour time jets of the
opposite-paired polynomial cutoffs and the inverse-square tail; it
does not assume absolute convergence of $P_1$ or arbitrary-sequence
well-posedness. Those hypotheses remain in force here.

\begin{theorem}[No nonconstant universal finite-head integral]
\label{thm:nohead}
Fix $P\ge1$. Let
\[
 I=I(t,D,S_1,\ldots,S_P,T_2,\ldots,T_P)
\]
be $C^1$ on an open real coordinate domain with $D\ne0$. Suppose
the chain-rule derivative of $I$ vanishes under
\eqref{eq:pair}, \eqref{eq:Shierarchy} and \eqref{eq:Thierarchy}
for freely variable tails through order $P+2$. Then $I$ is locally
constant. Here ``freely variable'' means that the relevant head and
tail coordinates range through an open set, with no extra relation
fixing the higher moments from the retained ones.
\end{theorem}
\begin{proof}
In $\dd I/\dd t$, the coefficients of $S_{P+2}$ and $T_{P+2}$
are respectively
\begin{equation}
 -4P(P+3)D I_{S_P},\qquad -4P(P+3)D I_{T_P}.
 \label{eq:highest-tail}
\end{equation}
The odd term is absent when $P=1$. No lower row contains either
highest coordinate, and the nonlinear terms contain neither.
Independent variation and $D\ne0$ imply that these partial
derivatives vanish. Locally remove those variables from $I$ and
repeat the argument at $P-1$. Eventually $I=I(t,D)$ and its
remaining equation is
\[
 I_t+(2-4DM)I_D=0.
\]
Variation of $M$ now forces $I_D=0$ and $I_t=0$. All first
derivatives of $I$ vanish on the domain, so it is locally constant.
\end{proof}

A $C^1$ first integral on an open neighborhood crossing $D=0$ cannot
acquire a nonconstant value only on that hypersurface: its derivatives
vanish on the dense complement and hence there too. This does not
exclude functions defined solely on the hypersurface, which is not an
invariant flow domain since $D'=2$ there.

\begin{proposition}[A finite polynomial realization of the free-tail premise]
\label{prop:rank}
At fixed real $D\ne0$, take $2P+3$ freely located real exterior
roots in a strict collar complement. There exist configurations where
the map from these roots to
\[
 (S_1,\ldots,S_{P+2},T_2,\ldots,T_{P+2})
\]
has full rank $2P+3$. Full rank persists on an open set and is generic
in any connected analytic configuration region on which the resulting
determinant is not identically zero.
\end{proposition}
\begin{proof}
The moment kernels are $f_r(w)=(w^2-D)^{-r}$ for $1\le r\le P+2$
and $g_r(w)=w(w^2-D)^{-r}$ for $2\le r\le P+2$. At infinity,
\[
 f_r'(w)=-2r w^{-2r-1}+O(w^{-2r-3}),\qquad
 g_r'(w)=(1-2r)w^{-2r}+O(w^{-2r-2}).
\]
Their leading powers are all distinct. In any putative nontrivial
linear relation, the least decaying nonzero term could not cancel.
They are consequently linearly independent rational functions.
Their evaluation vectors on an open real half-line outside the collar
span $\R^{2P+3}$: otherwise one nonzero fixed linear functional
would annihilate all evaluations, giving a rational identity there
and everywhere. Select distinct points whose evaluation vectors form
a basis. The moment Jacobian at those roots is precisely this
nonsingular evaluation matrix. Continuity and analyticity give the
stated open-set and generic conclusions.
\end{proof}

Each such initial divisor, together with the target pair, is a genuine
polynomial initial slice and has its finite heat evolution. Thus the
obstruction is not solely about fictitious independent sequences.
It remains a statement about a freely variable polynomial exterior
class. It does not establish free tails for the unchanged theta
source, exclude an invariant along one prescribed canonical
trajectory, treat an infinite functional, or cover a fixed small
divisor whose moments satisfy algebraic relations. In particular it
does not contradict Section~\ref{sec:finite}.

\section{Full-field constants and the Volterra boundary flux}
\label{sec:fullfield}

\begin{proposition}[Convergent heat pullback]
\label{prop:pullback}
For the canonical integral \eqref{eq:heat}, locally uniformly on
$\C^2$ in $(t,z)$,
\begin{equation}
 \mathcal J(t,z)=\sum_{r=0}^\infty\frac{t^r}{r!}
                       \partial_z^{2r}H(t,z)=H(0,z).
 \label{eq:pullback}
\end{equation}
Any fixed additional spatial derivative may be taken termwise.
\end{proposition}
\begin{proof}
After substituting the integral for $H$, the $r$th derivative adds
$(-u^2)^r$. The sum of the absolute integrands is bounded by
\[
 e^{(\Re t+|t|)u^2+|\Im z|u}\Phi(u).
\]
This is integrable uniformly on compact $(t,z)$ sets by the decay
of \eqref{eq:source}. A fixed extra derivative adds only a power
of $u$. Absolute interchange is valid, and the inner series is
$\sum_r(-tu^2)^r/r!=e^{-tu^2}$, which cancels the heat weight.
\end{proof}

For real $t$, $H(t,0)>0$. With $E_t(z)=H(t,z)/H(t,0)$,
\begin{equation}
 \frac{(e^{t\partial_z^2}E_t)(z)}{(e^{t\partial_z^2}E_t)(0)}
       =\frac{H(0,z)}{H(0,0)},\qquad
 (e^{t\partial_z^2}E_t)(0)=\frac{H(0,0)}{H(t,0)}>0.
 \label{eq:normalized-pullback}
\end{equation}
The parent's normalized even canonical product determines $E_t$
from the complete opposite-paired divisor. Thus
\eqref{eq:normalized-pullback} can encode the initial normalized
field using the full current divisor, without its amplitude.
It uses an infinite jet and explicit time. It does not give a
finite-head invariant or a paid truncation error, and it does not
establish any favorable zero property of $H(0,\cdot)$.

The positive-power sums of the full canonical divisor are not
absolutely convergent: the infinite roots do not tend to zero, so
even their terms in $\sum|\rho|^k$ fail the necessary condition
for $k>0$. Finite polynomial top-degree coefficients likewise have
no literal infinite-degree counterpart. Taking larger block or
product cutoffs in \eqref{eq:C4} requires a new normalization and
boundary-flux argument. Fixed-order convergence of the inverse
moments in \eqref{eq:moments} does not supply that argument.

\subsection{The Jacobi--Volterra realization}
The base manuscript \cite{BakerVolterra2026} derives the monic
orthogonal-polynomial recurrence for the even positive weight
$w_t(u)=e^{tu^2}\Phi(|u|)$ on $\R$ and the chain
\begin{equation}
 \beta_n'=\beta_n(\beta_{n+1}-\beta_{n-1}),\qquad
 \beta_0=0,\qquad \beta_n>0\ (n\ge1).
 \label{eq:volterra}
\end{equation}
Its finite prefix has a precise boundary flux:
\begin{equation}
 \frac{\dd}{\dd t}\sum_{n=1}^K\beta_n=\beta_K\beta_{K+1}.
 \label{eq:prefix}
\end{equation}
Indeed all internal products telescope in the sum of
\eqref{eq:volterra}. Closing a finite chain by setting
$\beta_{K+1}=0$ changes this physical law.

\begin{proposition}[Full support is not a finite trace invariant]
For each real $t$, the multiplication realization of the full Jacobi
operator has spectrum $\R$. Moreover $\sum_{n\ge1}\beta_n$ is not
finite. These statements require no bounded Lax propagator.
\end{proposition}
\begin{proof}
The weight is strictly positive and has full support. Multiplication
by $u$, on its natural domain in $L^2(w_t\dd u)$, is self-adjoint
with spectrum equal to the closed essential range $\R$. For example,
unit vectors supported in shrinking neighborhoods of any real
$\lambda$ are approximate eigenvectors, whereas nonreal $\lambda$
have bounded inverse multiplication by $(u-\lambda)^{-1}$.

For completeness, the polynomial realization is not missing an
orthogonal complement. If $f\in L^2(w_t\dd u)$ is orthogonal to all
polynomials, Cauchy--Schwarz and all exponential moments of $w_t$
make $\int f(u)w_t(u)e^{zu}\dd u$ entire. All derivatives at zero
vanish. Its Fourier transform is consequently zero; uniqueness for
the finite measure $f w_t\dd u$ gives $f=0$ almost everywhere.
Polynomials are dense. They are also a core for multiplication:
apply the same argument with weight $(1+u^2)w_t$ to its graph norm.
Thus the complete Jacobi realization is unitarily equivalent to
that unbounded multiplication operator.

If $\sum_n\beta_n<\infty$, then $\sup_n\beta_n<\infty$.
The Jacobi matrix with zero diagonal and off-diagonal
$\sqrt{\beta_n}$ would extend to a bounded operator, with norm at
most $2\sup_n\sqrt{\beta_n}$. Its closure could not equal the
unbounded multiplication operator just obtained. This contradiction
proves the assertion.
\end{proof}

The time-independent spectral set $\R$ is shared by many positive
full-support weights. It does not determine the zeros of their
Fourier readouts. Neither this spectral fact nor an unpaid infinite
trace substitutes for a bound on the Newman divisor.

\section{Finite-block invariants and complete exterior flux}
\label{sec:block}

Fix a real time interval and a stationary simple positively oriented
contour $\Gamma$ whose interior contains exactly $N$ zeros of a
holomorphic heat solution, counted with multiplicity. Assume $H$
does not vanish on a neighborhood of the contour throughout the
interval. Denote the interior divisor by $x_1,\ldots,x_N$ and set
\begin{equation}
 Q(t,z)=\prod_{i=1}^N(z-x_i(t)),\qquad
 R(t,z)=B(t,z)-\frac{Q_z(t,z)}{Q(t,z)}.
 \label{eq:exterior}
\end{equation}
The whole interior divisor has been removed, so $R$ is holomorphic
inside $\Gamma$. It includes the complete exterior and any analytic
factor of $H$; it is not a finite exterior cutoff. Assume also that
the interior contains the block's convex hull, so the center and all
Taylor segments used below lie inside. A cleared enclosing circle
is sufficient.

At a simple interior root the logarithmic derivative has regular
part $\sum_{j\ne i}(x_i-x_j)^{-1}+R(t,x_i)$, so
\begin{equation}
 x_i'=2\sum_{j\ne i}\frac1{x_i-x_j}+2R(t,x_i).
 \label{eq:block-roots}
\end{equation}
Let $c,y_i,p_k$ be the block quantities \eqref{eq:center} and
write $R_i=R(t,x_i)$. Pairwise cancellation, followed by subtraction
of the center velocity, gives
\begin{align}
 c'&=\frac2N\sum_iR_i, \label{eq:center-flux}\\
 p_2'&=2N(N-1)+4\sum_i y_iR_i,\notag\\
 p_3'&=6\sum_i(y_i^2-p_2/N)R_i,\notag\\
 p_4'&=4(2N-3)p_2+8\sum_i(y_i^3-p_3/N)R_i.
 \label{eq:block-powers}
\end{align}
For example, the exterior part of $p_k'$ is
$2k\sum_i(y_i^{k-1}-p_{k-1}/N)R_i$; the internal part is
\eqref{eq:powers}. This proves all lines of
\eqref{eq:block-powers}.

\subsection{Fourth order and the curvature remainder}
For $N\ge4$, retain the polynomial $C_4$ in \eqref{eq:C4} and put
\begin{equation}
 \beta=\frac{2N-3}{N(N-1)},\qquad
 \psi_i=y_i^3-\beta p_2y_i-p_3/N.
 \label{eq:psi}
\end{equation}
Differentiation with respect to a root gives
\begin{equation}
 \partial_{x_i}C_4=-\psi_i,\qquad
 \sum_i\psi_i=0,\qquad \sum_i\psi_iy_i=-4C_4.
 \label{eq:psi-identities}
\end{equation}
The first identity uses
$\partial_{x_i}p_k=k(y_i^{k-1}-p_{k-1}/N)$; the last follows by
substitution of \eqref{eq:C4}. Internal terms cancel in
\eqref{eq:block-powers}, giving the exact balance
\begin{equation}
 C_4'=-2\sum_i\psi_iR_i.
 \label{eq:C4flux}
\end{equation}
Now expand the complete exterior at the current mean:
\begin{equation}
 R(t,c+y)=a(t)+b(t)y+r(t,y),\quad
 a=R(t,c),\quad b=R_z(t,c).
 \label{eq:affine}
\end{equation}
The identities \eqref{eq:psi-identities} yield
\begin{equation}
 C_4'-8bC_4=\E_4,\qquad
 \E_4=-2\sum_i\psi_i r(t,y_i).
 \label{eq:affine-C4}
\end{equation}
In particular an exactly affine exterior gives the constant
$e^{-8\int_{t_0}^t b(v)\dd v}C_4(t)$. For a general exterior,
the same multiplier $\mu(t)$ gives the exact paid balance
\begin{equation}
 \mu(t)C_4(t)-C_4(t_0)=\int_{t_0}^t\mu(v)\E_4(v)\dd v.
 \label{eq:integrated-C4}
\end{equation}

\begin{proposition}[A local curvature allowance]
\label{prop:curvature}
Suppose $|y_i|\le\rho$ and $|R_{zz}|\le C_2$ on every segment
from $c$ to $x_i$. Then
\begin{equation}
 |\E_4|\le C_2\sum_i|\psi_i||y_i|^2\le4NC_2\rho^5.
 \label{eq:curvature}
\end{equation}
\end{proposition}
\begin{proof}
Taylor's integral remainder gives
$r(t,y)=y^2\int_0^1(1-s)R_{zz}(t,c+sy)\dd s$, hence
$|r(t,y)|\le C_2|y|^2/2$. Also
\[
 |\psi_i|\le\rho^3+\beta N\rho^3+\rho^3
          =\left(4-\frac1{N-1}\right)\rho^3\le4\rho^3.
\]
Apply these bounds to \eqref{eq:affine-C4}.
\end{proof}

The fifth power is a local scale allowance. A near-conservation
estimate over a time interval still needs bounds on the initial
value, $b$, $C_2$ and $\rho$ throughout that interval.

\subsection{The general homogeneous law}
\begin{theorem}[Affine-corrected flux for a homogeneous heat integral]
\label{thm:general-flux}
Let $J(x_1,\ldots,x_N)$ be a symmetric polynomial, invariant under
common translation and homogeneous of degree $m\ge1$. Assume it
is an autonomous first integral of the isolated finite heat-root
law \eqref{eq:finite-roots}. Write $J_i=\partial_{x_i}J$ and use
the complete exterior \eqref{eq:exterior}. Then
\begin{align}
 J'&=2\sum_i J_iR_i, \label{eq:Jflux}\\
 J'-2mbJ&=\E_J=2\sum_iJ_i r(t,y_i). \label{eq:Jaffine}
\end{align}
If $Z=z-c$, the single rational kernel
\begin{equation}
 K_J(z)=\sum_i\frac{J_i}{z-x_i}-\frac{mJ}{(z-c)^2}
       =\frac1{Z^2}\sum_i\frac{J_i y_i^2}{Z-y_i}
 \label{eq:Kgeneral}
\end{equation}
satisfies
\begin{equation}
 \E_J=2\av{B(t,z)K_J(z)}.
 \label{eq:contour-general}
\end{equation}
All these symmetric statements extend through internal mergers
while the count and outer clearance persist. For the integrals
$J_j$ of Corollary~\ref{cor:homogeneous}, take $m=j$.
\end{theorem}
\begin{proof}
The chain rule and the assumed isolated first-integral identity
cancel the internal part of \eqref{eq:block-roots}, leaving
\eqref{eq:Jflux}. Translation invariance and Euler's homogeneous
identity give
\begin{equation}
 \sum_iJ_i=0,\qquad \sum_i y_iJ_i=mJ.
 \label{eq:Euler}
\end{equation}
Substitution of \eqref{eq:affine} into \eqref{eq:Jflux} now proves
\eqref{eq:Jaffine}. In particular a purely affine exterior gives
the constant $e^{-2m\int b}J$.

For $w$ inside the contour, Cauchy's formula and the complete
subtraction in \eqref{eq:exterior} give
\begin{equation}
 R(t,w)=\av{\frac{B(t,z)}{z-w}},\qquad
 R_z(t,c)=\av{\frac{B(t,z)}{(z-c)^2}}.
 \label{eq:exterior-Cauchy}
\end{equation}
To verify that the interior poles do not add a term, integrate
$(Q_z/Q)/(z-w)$: all its poles are inside, and it is $O(z^{-2})$
at infinity, so their total residue is zero. The derivative
formula follows by differentiating with respect to $w$. These
statements remain true when $w$ coincides with a block root.
Substitution in $2\sum_iJ_iR_i-2mbJ$ gives
\eqref{eq:contour-general}. Finally,
\[
 \frac1{Z-y_i}=\frac1Z+\frac{y_i}{Z^2}
                   +\frac{y_i^2}{Z^2(Z-y_i)}
\]
and \eqref{eq:Euler} prove the second equality in
\eqref{eq:Kgeneral}. The $Z^{-1}$ and $Z^{-2}$ terms cancel exactly.

For continuation, each raw power sum of the block is
$\av{z^kB}$ and is holomorphic in time on a cleared contour.
Newton identities recover the coefficients of $Q$, and hence every
symmetric polynomial used above. The quotient $H/Q$ is nonvanishing
inside, and its logarithmic derivative $R$ is holomorphic there
and in time. Expressing $J$ as a polynomial in the raw power sums
shows by differentiation that $J_i$ is a polynomial in $x_i$ with
symmetric coefficient functions;
therefore $\sum_iJ_iR(x_i)$ is a finite combination of holomorphic
counted traces. Equivalently, at a multiple slice perturb the heat
solution to $H+\epsilon$, which has simple block zeros for generic
small $\epsilon$, preserves the boundary count and clearance, and
converges with its derivatives there. All displayed symmetric
quantities and first time derivatives converge by their contour
formulas. Passing to the limit proves the identities at the merger,
without assigning a singular individual velocity at that time.
\end{proof}

For $J=C_4$, $J_i=-\psi_i$ and $m=4$. Thus our sign convention
is $K_{C_4}=-\sum_i\psi_i/(z-x_i)-4C_4/(z-c)^2$; it is the
negative of the founding note's $K_4$. The flux
$2\av{BK_{C_4}}$ is exactly the same $\E_4$.
The separate curvature channel is also available from
\begin{equation}
 R_{zz}(t,w)=2\av{\frac{B(t,z)}{(z-w)^3}}.
 \label{eq:Rzz}
\end{equation}
On a circle of radius $A_0$, a uniform source error $\epsilon$
and a distance $\delta$ from $w$ to that circle contribute at most
$2A_0\epsilon/\delta^3$ to this derivative. The next section
retains the cancellations in the combined kernel before taking
absolute values.

\section{One contour error comparison, with coordinate costs}
\label{sec:errors}

The bounded follow-up fixes one source-error model: a uniform
pointwise error in $B$ on one cleared circle. It compares two
explicit triangle/Cauchy allowances using the same data. It does
not compare exact feasible sets, source-acquisition costs or
canonical minimum-gap widths. Direct contour evaluation is
allowed to use the same combined kernel.

\begin{theorem}[Quadratic cancellation factor]
\label{thm:quadratic}
Use the exact block coordinates of Theorem~\ref{thm:general-flux}.
Let $\Gamma$ be the circle $|z-c|=A$, where
$|y_i|\le\rho<A$. Suppose $\widehat B$ satisfies
$\sup_\Gamma|B-\widehat B|\le\epsilon$. For the linear functional
$\E_J=2\av{BK_J}$, separate evaluation of the exterior values
and $b$ gives the valid allowance
\begin{equation}
 B_{\rm sep}=2A\epsilon\sum_i\frac{|J_i|}{A-|y_i|}
                      +\frac{2m|J|\epsilon}{A}.
 \label{eq:Bsep}
\end{equation}
The combined kernel gives the valid allowance
\begin{equation}
 B_{\rm comb}=\frac{2\epsilon}{A}
                  \sum_i\frac{|J_i||y_i|^2}{A-|y_i|},
 \qquad
 B_{\rm comb}\le\left(\frac\rho A\right)^2 B_{\rm sep}.
 \label{eq:Bcomb}
\end{equation}
The inequality is meaningful without dividing by $B_{\rm sep}$,
including cases where it is zero.
\end{theorem}
\begin{proof}
The contour length is $2\pi A$. From
\eqref{eq:exterior-Cauchy}, the error in each $R_i$ is bounded by
$A\epsilon/(A-|y_i|)$ and the error in $b$ by $\epsilon/A$.
Applying the triangle inequality to
$2\sum_iJ_iR_i-2mJb$ proves \eqref{eq:Bsep}.
The second form of \eqref{eq:Kgeneral} gives
\[
 \sup_\Gamma|K_J|\le\frac1{A^2}
                   \sum_i\frac{|J_i||y_i|^2}{A-|y_i|}.
\]
Multiply by $2A\epsilon$ to obtain $B_{\rm comb}$. Replacing
$|y_i|^2$ by $\rho^2$ proves
\[
 B_{\rm comb}\le(\rho/A)^2
       \left(2A\epsilon\sum_i\frac{|J_i|}{A-|y_i|}\right)
       \le(\rho/A)^2 B_{\rm sep}.
\]
\end{proof}

This is a strict factor less than one whenever the separate
allowance is positive and $\rho<A$. It is a comparison with the
specific allowance \eqref{eq:Bsep}; the latter is not asserted
to be the sharp independent-moment baseline. No superiority over
a direct method using \eqref{eq:Kgeneral} follows. This distinction
matches the parent's common-data comparison discipline
\cite{BakerCorrelation2026}.

\subsection{A quartet refinement and its asymmetry term}
\begin{proposition}[The next cancellation for a symmetric quartet]
\label{prop:cubic}
For $N=4$ and $J=C_4$, let $q_2=\sum_iJ_i y_i^2$. Then
\begin{equation}
 q_2=-\frac16p_2p_3,\qquad
 K_{C_4}(z)=\frac{q_2}{Z^3}
          +\frac1{Z^3}\sum_i\frac{J_i y_i^3}{Z-y_i}.
 \label{eq:cubic-kernel}
\end{equation}
Under Theorem~\ref{thm:quadratic}'s assumptions another valid
source-error allowance is
\begin{equation}
 B_{\rm cubic+asym}=\frac{2\epsilon|q_2|}{A^2}
       +\frac{2\epsilon}{A^2}
             \sum_i\frac{|J_i||y_i|^3}{A-|y_i|}.
 \label{eq:Bcubic}
\end{equation}
If the centered quartet is reflection-symmetric, $p_3=q_2=0$,
and
\begin{equation}
 B_{\rm cubic}\le(\rho/A)^3 B_{\rm sep}.
 \label{eq:cubic-factor}
\end{equation}
No parity condition is imposed on the source error $B-\widehat B$.
\end{proposition}
\begin{proof}
Since $J_i=-y_i^3+\beta p_2y_i+p_3/N$,
$q_2=-p_5+(\beta+1/N)p_2p_3$. For a centered quartet,
Newton identities give
$p_5=-b_2p_3-b_3p_2=5p_2p_3/6$, while
$\beta+1/N=5/12+1/4=2/3$. This proves the first identity.
Expanding one more term in \eqref{eq:Kgeneral} proves the second.
The circle-length bound gives \eqref{eq:Bcubic}. When $p_3=0$,
replace $|y_i|^3$ by $\rho^3$ and compare with the first term of
\eqref{eq:Bsep}, as in Theorem~\ref{thm:quadratic}.
\end{proof}

Approximate symmetry alone cannot set $q_2$ to zero. The first
term of \eqref{eq:Bcubic} is its explicit payment; even that
allowance need not improve \eqref{eq:Bcomb} in every asymmetric
configuration.

\subsection{Uncertain block coordinates}
The preceding factors keep the coordinates exact. A root enclosure
must pay for the kernel's coordinate dependence as well. The
following conservative bound supplies one such payment for $C_4$.

\begin{proposition}[A coordinate allowance for the fourth-order kernel]
\label{prop:coordinates}
Let $x$ and $\widehat x$ be two labeled $N$-tuples, $N\ge4$, with
$\max_i|x_i-\widehat x_i|\le e$. On the straight homotopy between
them, assume every centered root has modulus at most $\rho$.
Let one fixed circle $\Gamma$ have radius $A_0$. Suppose for every
homotopy parameter and every $z\in\Gamma$ that
\begin{equation}
 |z-x_i|\ge\delta>0,\qquad |z-c|\ge\delta>0.
 \label{eq:coordinate-clearance}
\end{equation}
For $K=K_{C_4}$,
\begin{equation}
 \sup_\Gamma|K(x)-K(\widehat x)|\le G,
 \quad
 G=Ne\left(\frac{20\rho^2}{\delta}
          +\frac{20\rho^3}{\delta^2}
          +\frac{6\rho^4}{\delta^3}\right).
 \label{eq:G}
\end{equation}
If $|\widehat B|\le B_*$ and $|B-\widehat B|\le\epsilon$ on
that circle, the corresponding exact contour functionals obey
\begin{equation}
 |\E_4(B,x)-\E_4(\widehat B,\widehat x)|
 \le 2A_0\left[\epsilon\sup_\Gamma|K(\widehat x)|
                         +(B_*+\epsilon)G\right].
 \label{eq:joint-coordinate}
\end{equation}
Here $\E_4(B,x)$ denotes the functional in
\eqref{eq:contour-general}; its identification with the physical
flux uses the actual counted block $x$.
\end{proposition}
\begin{proof}
Differentiating \eqref{eq:psi} in a coordinate gives
\begin{equation}
 \partial_{x_j}\psi_i
  =(3y_i^2-\beta p_2)(\delta_{ij}-1/N)
       -2\beta y_i y_j-\frac3N(y_j^2-p_2/N).
 \label{eq:psi-jacobian}
\end{equation}
Since $\beta N<2$, the sums of absolute values of these three
terms over $j$ are at most $10\rho^2$, $4\rho^2$ and
$6\rho^2$, respectively. Thus
$\sum_j|\partial_{x_j}\psi_i|\le20\rho^2$.
Also $|J_i|=|\psi_i|\le4\rho^3$ and, from \eqref{eq:C4},
$|C_4|\le3N\rho^4/4$. Along the homotopy, these imply
\[
 |\dot J_i|\le20\rho^2 e,\qquad
 |\dot C_4|\le4N\rho^3 e,\qquad |\dot c|\le e.
\]
Differentiate
$K=\sum_iJ_i/(z-x_i)-4C_4/(z-c)^2$ along the homotopy:
\[
 \dot K=\sum_i\left(\frac{\dot J_i}{z-x_i}
                       +\frac{J_i\dot x_i}{(z-x_i)^2}\right)
       -\frac{4\dot C_4}{(z-c)^2}
       -\frac{8C_4\dot c}{(z-c)^3}.
\]
The four contributions are bounded by
$20N\rho^2e/\delta$, $4N\rho^3e/\delta^2$,
$16N\rho^3e/\delta^2$ and $6N\rho^4e/\delta^3$.
Integrate over the unit homotopy interval to obtain \eqref{eq:G}.
Finally use the exact splitting
\[
 BK(x)-\widehat B K(\widehat x)
  =(B-\widehat B)K(\widehat x)+B\bigl(K(x)-K(\widehat x)\bigr)
\]
and $|B|\le B_*+\epsilon$ in the circle-length bound. This also
pays the mixed field/coordinate error.
\end{proof}

If the fixed circle is centered at the nominal mean, the cancellation
bounds apply to $K(\widehat x)$. Its distance to the actual mean
and roots must still satisfy \eqref{eq:coordinate-clearance} along
the whole homotopy. The term involving $G$ can erase the nominal
source-error gain. A numerical contour quadrature introduces a
separate error for the integrand actually integrated; arithmetic
enclosures and any integration of \eqref{eq:integrated-C4} in time
must be paid too. Counts, clearance, and between-time control are
additional premises, not consequences of the radius factor.

\subsection{Two fixed exact arithmetic checks}
The bounded check uses two six-root polynomial initial slices,
and no canonical source acquisition. For the block
$(-2,-1,1,2)$ with exterior $(-5,5)$, exact rational evaluation gives
\begin{equation}
 C_4=\frac{23}{12},\quad C_4'=-\frac{55}{42},\quad
 b=-\frac2{25},\quad \E_4=-\frac{29}{350}.
 \label{eq:symmetric-example}
\end{equation}
The same block without the exterior has $C_4'=0$. At radius
$A=4$ and unit source-error scale, the specified allowances are
\begin{equation}
 B_{\rm sep}=\frac{421}{18},\quad
 B_{\rm comb}=\frac{31}{18},\quad
 B_{\rm cubic}=\frac{43}{72}.
 \label{eq:example-allowances}
\end{equation}
Thus their exact ratios to the separate allowance are $31/421$
and $43/1684$. These are allowance ratios on exact finite geometry,
not measured errors or canonical width reductions.

For the asymmetric block $(-3,-1,2,5)$ with exterior $(-8,9)$,
the center is $3/4$. At radius $A=6$ the roots have
$\rho=17/4$, and exact arithmetic gives
\begin{align}
 p_2&=147/4,&p_3&=165/8,&C_4&=861/128,\notag\\
 q_2&=-8085/64,&C_4'&=-6847/7280,&
 \E_4&=7306487/13213200.
 \label{eq:asymmetric-example}
\end{align}
This checks the nonzero asymmetry payment. The execution record
\cite{BakerAffineChecks2026} contains both sets of allowance values
and a quartic collision at $T=1/20$, $b=2$. In each six-root case
it computes the derivative from all root velocities, the exterior
balance, and the exterior residues of the combined contour kernel
independently in exact rational arithmetic. These checks confirm
the finite algebra; the general theorems rest on their proofs.

\section{Collision controls, evidence, and the remaining question}
\label{sec:controls}

The conservation statements are generic heat identities. They cannot
distinguish the unchanged canonical theta source from another source
merely because that source also obeys the heat equation. In particular,
the positive value in \eqref{eq:positive-collision} coexists with an
exact double event. The finite invariant plus an independently paid
geometric condition can be informative, as
\eqref{eq:quartic-discriminant} shows, but its sign alone is not
a collision criterion.

The following analytical premise screen uses the documented witnesses
in the project's collision catalog \cite{ProjectControls2026}.
Only the quartic row is newly checked by the present bounded program.
The other rows do not claim new numerical validation of those witnesses.
None of the generic identities rejects a genuine heat-flow control
at its collision. A future exclusion estimate must fail there, or
must identify an additional premise that the control does not meet.

\begingroup
\small
\setlength{\tabcolsep}{5pt}
\begin{longtable}{@{}p{.23\textwidth}p{.70\textwidth}@{}}
\toprule
Control class & Conservation outcome and missing canonical premise\\
\midrule
\endfirsthead
\toprule
Control class & Conservation outcome and missing canonical premise\\
\midrule
\endhead
\bottomrule
\endfoot
Shifted quartic; isolated quadratic &
The finite coefficient integrals survive the multiple event. The
quartic explicitly defeats positivity of $C_4$ as a barrier. These
finite polynomials do not have the literal theta source or infinite
canonical population.\\[5pt]
Cosine plus a constant &
Count conservation and its own heat pullback survive the documented
double zero. Its source is not the complete theta source; an identity
does not supply a favorable canonical flux.\\[5pt]
Euler-factor translates &
The complete block balances apply to the changed field. Unchanged
normalized theta coefficients and completion fail; the original
source's quantitative flux cannot be transferred without proof.\\[5pt]
Weighted-square sources &
Square support and selected endpoint equalities coexist with the
documented collision. Full reflected, normalized unit-weight theta
source equality fails. The generic conservation identities remain.\\[5pt]
Reflected super-Gaussian &
Local heat and block laws still apply. Literal theta arithmetic and
the canonical order-one hypothesis fail, so the global canonical
product step is not imported into this control.\\[5pt]
Smooth shape; positive primitive &
Positivity, evenness, strong decay and the listed auxiliary endpoint
or primitive properties do not fix the exterior flux. Complete theta
arithmetic/profile equality fails. The identities do not reject the
catalog's collision constructions.\\[5pt]
Finite-prime, finite-score and finite-profile variants &
Preserving a finite set of observations does not provide unchanged
full source equality or a paid joint flux. Their particular absent
reflection, arithmetic or full-profile premise remains the one
specified for each witness in the catalog.\\[5pt]
Changed observation kernels; fractional divisor profiles &
The full cosine/Newman equation, or respectively an integer divisor
realization, is missing. They are information-sufficiency controls,
not counterexamples to a theorem about the physical heat divisor.\\
\end{longtable}
\endgroup

\subsection{What has and has not been obtained}
The finite coefficient integrals, the specialized nilpotence, the
free-tail obstruction, the full-field pullback, and the fourth-order
block balance have been derived from the founding exploration.
The bounded extension proves a general homogeneous exterior-flux
identity, its merger continuation, the specified quadratic source-error
factor, the quartet refinement and an explicit coordinate payment.
The checks use exact rational arithmetic on two fixed six-root
configurations and one quartic event. The inherited finite check
record also tests polynomial pullback, autonomous integrals,
nilpotence and examples of moment-map rank
\cite{BakerConservationBrainstorm2026}.

The evidence level is \emph{proved in this working draft, with exact
finite algebra checks and same-assistant assessment}. This is not an
independent or specialist review. The external reading scope is also
limited: \cite{DeconinckKimuraSegur2007} was inspected at metadata and
abstract level, while slide 4 and slides 27--28 of
\cite{Sokal2023HeatZeros} were checked. Their finite-integrability
context is credited, but their uninspected arguments and conjectures
are not used as proof premises. The repository reference catalog
retains this provenance and the failed full-text retrieval.

The source-error reduction is proved only against the displayed
separate Cauchy allowance with the same uniform field error and
exact block geometry. Direct evaluation can retain the same kernel.
The coordinate term in \eqref{eq:joint-coordinate}, contour
quadrature, count and clearance, and whole-time integration remain
part of any total comparison. No canonical block has been acquired
or evaluated in this branch, and no total canonical gap-width or
cost improvement is claimed. No RH or Newman-threshold consequence
is established.

The useful remaining question is therefore concrete: on one already
counted canonical block and one fixed physical interval, can the
complete balance and its shared contour cancellations improve a
joint statistic or gap enclosure after every source, coordinate,
quadrature, arithmetic and time cost is paid? The comparison must
give a direct centered-contour method the same information and
the same available cancellations. Answering that question requires
a separately selected source record and consumer. It is not
activated by the present analytical draft, and it does not change
the parent investigation's current source or two-sided-refinement
work.

\bibliographystyle{plainurl}
\bibliography{branch,references}
\end{document}
